\documentclass[11pt]{article}
\usepackage[utf8]{inputenc}
\usepackage[margin=0.82in]{geometry}
\setlength{\parindent}{0pt}
\setlength{\parskip}{0.68em}
\linespread{1.00}
\pagestyle{empty}

\begin{document}

Dear Editors of the Journal of Regional Science,

I am pleased to submit my manuscript, ``Public Co-Creation Hubs, Cross-Regional Pipelines, and Dynamic Efficiency in Thin Regional Ecosystems,'' for consideration for publication in the Journal of Regional Science.

This paper develops a two-region, two-period model of public co-creation hubs in thin regional ecosystems. The paper is motivated by a simple but increasingly important regional policy question: when does a publicly supported hub genuinely improve welfare in a thin regional ecosystem, and when should it instead be evaluated as one component of a broader regional strategy that also creates external renewal? In the model, a public hub raises local matching within a target region, while a complementary pipeline policy creates cross-regional renewal through outside linkages.

The analysis yields three main results. First, a hub is socially efficient as a stand-alone intervention only when its local coordination gains are sufficiently large to offset both fixed cost and the dynamic redundancy generated by repeated local interaction in a thin regional pool. Second, visible short-run success and social efficiency are not equivalent: a hub may perform well under first-period activity metrics while remaining socially undesirable once novelty loss and lock-in risk are taken into account. Third, the relevant policy object is often not a hub alone, but a bundled regional strategy combining local coordination with cross-regional pipeline formation.

I believe the manuscript is particularly well suited to the Journal of Regional Science because it speaks directly to central concerns of regional science: regional innovation, place-based development policy, local buzz, external linkages, and the conditions under which regional coordination mechanisms improve welfare. More specifically, the paper connects agglomeration and proximity arguments to an explicit dynamic welfare framework, and it shows why policies that appear successful on visible local activity measures may nevertheless be inefficient unless they also sustain cross-regional renewal. In that sense, the paper contributes not only to the theory of regional coordination, but also to the evaluation of contemporary regional hub policies.

The paper also generates empirically testable implications. It predicts that hub creation should raise initial local interaction volume, but not necessarily sustain the share of new ties; that hub-plus-pipeline regions should exhibit stronger outside collaboration and higher medium-run novelty than hub-only regions; and that the welfare value of hub policy should be greater in regions with stronger absorptive capacity. These implications, I believe, further strengthen the paper's relevance for the Journal of Regional Science readership.

This manuscript is original, has not been published previously, and is not under consideration elsewhere. Thank you for your consideration.

\vfill

Sincerely,

Ryota Matsuki\\
Independent Researcher\\
Matsuyama, Ehime, Japan\\
ryota.matsuki@gmail.com

\end{document}
